\documentclass[11pt]{article}
\usepackage[margin=0.9in]{geometry}
\usepackage[T1]{fontenc}
\usepackage{lmodern,amsmath,amssymb,amsthm,microtype}
\usepackage[colorlinks=true,urlcolor=blue,citecolor=blue,linkcolor=blue]{hyperref}
\theoremstyle{definition}
\newtheorem{theorem}{Theorem}
\newtheorem{lemma}[theorem]{Lemma}
\newtheorem{corollary}[theorem]{Corollary}
\newcommand{\Sym}{\operatorname{Sym}}
\newcommand{\WLP}{\mathrm{WLP}}
\newcommand{\SLP}{\mathrm{SLP}}
\title{A symmetric ideal chain with periodic Lefschetz\\ behavior in positive characteristic}
\author{Alec Kriebel\\ \small Independent researcher\\
\small \href{https://orcid.org/0009-0001-9320-500X}{ORCID 0009-0001-9320-500X}}
\date{22 October 2026 \quad Version 1.0.0}
\begin{document}
\maketitle
\begin{abstract}
For each fixed field of positive characteristic $p$, we construct a single
compatible symmetric chain of Artinian ideals in two rows of variables, generated
from width three. At every width $n\geq2$, its quotient has both the weak and
strong Lefschetz properties exactly when $p\nmid n-1$. The quadratic quotient is
free over the integers, and a determinant for an arbitrary linear form gives the
classification over every field extension. In particular, over the fixed infinite
field $\mathbb F_p(t)$ both properties hold and fail at unbounded widths.
This gives a negative answer to the unrestricted eventual dichotomies posed by
R\"omer. It gives no negative answer in characteristic zero or in one row.
\end{abstract}

\section{The question and the chain}
Let $K$ be a fixed field and $c\geq1$ a fixed integer. In the setup of
R\"omer~\cite[pp.~1594--1596]{romer},
$R_n=K[X_{i,j}:1\leq i\leq c,\ 1\leq j\leq n]$ and $\Sym(n)$
permutes the column index $j$, simultaneously in every row.
A symmetric invariant chain consists of ideals $I_n\subseteq R_n$ such that
each $I_n$ is $\Sym(n)$-invariant and
\[
 \Sym(n+1)\bigl(I_nR_{n+1}\bigr)\subseteq I_{n+1}.
\]
For Artinian homogeneous quotients, R\"omer asks whether each Lefschetz property
eventually always holds or eventually always fails. His contribution fixes an
arbitrary field and $c\geq1$; the separate one-row monomial working-group
discussion on pp.~1544--1545 is a special setting.

We use the usual definitions: a standard graded Artinian $K$-algebra $A$
has $\WLP$ if some $\ell\in A_1$ makes every map
$\ell:A_d\to A_{d+1}$ have rank $\min(\dim A_d,\dim A_{d+1})$.
It has $\SLP$ if one $\ell$ has this property for every
$\ell^s:A_d\to A_{d+s}$, $d\geq0$, $s\geq1$.
These are existence definitions over $K$. Over an infinite field the
successful forms form a nonempty Zariski-open set, when they exist,
since the conditions are finitely many nonvanishing minors. Standard
background on characteristic-sensitive Lefschetz maps is given in~\cite{cook}.
Our successful forms are already defined over the prime field.

Set $c=2$, write the variables as $x_1,\ldots,x_n,y_1,\ldots,y_n$, and put
$\mathcal m_n=(x_1,\ldots,x_n,y_1,\ldots,y_n)$. Define
\begin{equation}\label{ideal}
 I_n=\mathcal m_n^3+
 (x_ix_j:i\ne j)+(y_iy_j:1\leq i,j\leq n)+(x_iy_i:1\leq i\leq n)
 +(x_iy_j-x_iy_k:i,j,k\text{ pairwise distinct}).
\end{equation}
Here $\mathcal m_n^3$ means all cubic monomials; in particular repeated
variables are included. Empty index conditions contribute no generators.
All generators are homogeneous and the entire list is preserved by column
permutations. Every injection of column sets carries a displayed generator
to one of the same type. This proves the chain compatibility, including
the first two widths. Since $\mathcal m_n^3\subseteq I_n$, every
$A_n=R_n/I_n$ is Artinian.

For $n\geq3$, in fact $I_n=(\Sym(n)I_3)R_n$.
One containment follows from compatibility. For the other, every generator
in~\eqref{ideal} involves at most three columns, including every cubic.
Inject those distinct supporting columns into $\{1,2,3\}$, without changing
which indices repeat. The resulting generator belongs to $I_3$.
The reverse injection on its support extends to a permutation of
$\{1,\ldots,n\}$ by bijecting the unused columns. Thus the generator is a
translate of one at width three. This uses one fixed finite list of orbits,
rather than choosing ideals according to the residue of $n$.

\begin{theorem}\label{main}
For every field $K$, the quotients in~\eqref{ideal} satisfy the following.
At $n=1$, $A_1$ has $\SLP$. At every $n\geq2$, its Hilbert function is
$(1,2n,2n)$, and
\[
 A_n\text{ has }\SLP
 \ \Longleftrightarrow\ A_n\text{ has }\WLP
 \ \Longleftrightarrow\ (n-1)\,1_K\ne0.
\]
These assertions remain true after any field extension.
\end{theorem}

\section{Integral basis and multiplication}
The same generators define an ideal over $\mathbb Z$. In degree two first
remove the monomials $x_ix_j$ with $i\ne j$, all $y_iy_j$, and all $x_iy_i$.
For $n\geq2$ the remaining symbols are $x_i^2$ and $x_iy_j$ with $j\ne i$.
Fix $i$ and choose one reference $j_i\ne i$. The $n-1$ symbols
$x_iy_j$ have relations precisely all their pairwise differences.
The reference symbol, together with
$x_iy_j-x_iy_{j_i}$ for $j\ne i,j_i$, is a unimodular basis for their
free abelian group. The relation subgroup is exactly the span of those
differences. Its quotient is therefore free of rank one, with basis
$v_i=[x_iy_{j_i}]$. When $n=2$ there is one symbol and no differences.
Different $i$ do not interact, and the $u_i=[x_i^2]$ are unrestricted.
Consequently the degree-two quotient has the integral basis
\[
 u_1,\ldots,u_n,v_1,\ldots,v_n.
\]
There are no linear relations, and all degrees at least three vanish.
Thus the whole integral graded quotient is torsion free and has graded
ranks $(1,2n,2n)$. Base change to any field gives exactly the claimed
basis and Hilbert function, without a characteristic-dependent dimension
argument.

Fix $n\geq2$ and an arbitrary linear form
\[
 \ell=\sum_{i=1}^n(a_ix_i+b_iy_i),\qquad B=\sum_{i=1}^n b_i.
\]
The defining products give, for each $j$,
\[
 \ell x_j=a_j u_j+(B-b_j)v_j,\qquad
 \ell y_j=\sum_{i\ne j}a_i v_i.
\]
Let $D(a)$ and $D(B-b)$ denote the corresponding diagonal matrices and
let $J$ be the $n$ by $n$ all-ones matrix. In the ordered source basis
$(x_1,\ldots,x_n,y_1,\ldots,y_n)$ and target basis $(u_1,\ldots,u_n,
v_1,\ldots,v_n)$, multiplication is
\begin{equation}\label{matrix}
 M_\ell=
 \begin{pmatrix}
 D(a)&0\\
 D(B-b)&D(a)(J-I)
 \end{pmatrix},\qquad
 \det M_\ell=(-1)^{n-1}(n-1)\prod_{i=1}^n a_i^2 .
\end{equation}
For completeness, expand $\det(J-I)$ multilinearly in its columns over
$\mathbb Z$. Choosing no column from $J$ contributes $(-1)^n$.
Choosing exactly one contributes $(-1)^{n-1}$ for each of the $n$ choices.
Choosing at least two gives equal columns and determinant zero.
Their sum is $(-1)^{n-1}(n-1)$. This proves~\eqref{matrix}
over every field, without division by $n$ or an eigenbasis assumption.

\section{Bad widths, good widths, and the initial width}
If $(n-1)1_K=0$, equation~\eqref{matrix} is zero for \emph{every}
choice of coefficients. Since $(A_n)_1$ and $(A_n)_2$ both have dimension
$2n$, no multiplication between them can have maximal rank.
There is also a direct obstruction independent of the determinant:
\[
 z=\sum_{j=1}^n y_j\ne0,\qquad
 y_i z=0,\qquad x_i z=(n-1)v_i=0.
\]
The nonzero element $z$ is in the degree-one socle, so every linear form
annihilates it. Its nonvanishing follows from the degree-one basis,
regardless of whether the characteristic divides $n$. This obstruction
also survives every field extension. In particular, failure is not inferred
merely from testing one invariant form.

If $(n-1)1_K\ne0$, take $\ell=x_1+\cdots+x_n$.
Then $M_\ell=\operatorname{diag}(I,J-I)$ is invertible.
The map $(A_n)_0\to(A_n)_1$ is injective, and maps with zero target
have maximal rank. Thus $\ell$ is weak Lefschetz.
Furthermore,
\[
 \ell^2=\sum_{i=1}^n u_i\ne0,
\]
since all off-diagonal $x_ix_j$ vanish and the $u_i$ are independent.
This remains valid in characteristic two; it does not rely on a
nonzero cross-term coefficient. The only additional map with nonzero
source and target needed for $\SLP$ is
$\ell^2:(A_n)_0\to(A_n)_2$, which is injective.
All other positive powers either were checked or have zero target.
Thus $\ell$ is strong Lefschetz. Conversely, $\SLP$ includes the
maps for $s=1$ and implies $\WLP$.

At $n=1$, the quotient is
$K[x,y]/(xy,y^2,(x,y)^3)$, with basis $1,x,y,x^2$ and Hilbert
function $(1,2,1)$. Multiplication by $x$ from degree zero is injective,
from degree one to two is surjective, and multiplication by $x^2$
from degree zero to two is nonzero. All remaining positive-power maps
have zero target. Hence $x$ is strong Lefschetz. This completes the
proof of Theorem~\ref{main}.

\begin{corollary}
For each prime $p$, over the single fixed infinite field $K=\mathbb F_p(t)$,
the chain~\eqref{ideal} has both $\WLP$ and $\SLP$ at every width
$n=2+kp$ with $k\geq0$, and neither property at every width
$n=1+kp$ with $k\geq1$. Consequently neither unrestricted eventual
dichotomy in R\"omer's question holds.
\end{corollary}
\begin{proof}
The theorem applies to all widths of this one chain over this one field.
Both displayed sequences are unbounded. Thus neither eventual constant
truth value is possible for either property.
\end{proof}
For $n\geq2$ the full pattern has period $p$: only $n\equiv1\pmod p$
is bad. In characteristic two, even widths are good and odd widths
at least three are bad. In characteristic zero this construction is strong
Lefschetz at every width. No conclusion about all characteristic-zero
chains or the one-row question follows. If desired, for a fixed $c>2$
one can add every variable in each extra row as a linear generator;
the quotient is the same. The two-row example alone disproves the
formulation allowing arbitrary fixed $c\geq1$.

\section{Related work, provenance, and verification limits}
Seceleanu and \c Sega's \emph{General symmetric ideals}~\cite{ss}
studies general ideals in a single row. Its Theorem B proves $\WLP$
under the hypothesis $\operatorname{char}K=0$ or
$\operatorname{char}K>n$, together with its stated size conditions.
Its chain result, Theorem 9.2 (summarized in Theorem C), assumes
characteristic zero. A fixed positive characteristic cannot remain larger
than unbounded $n$. These results therefore do not give the two-row
fixed-characteristic construction above. The underlying use of ranks,
socles, and characteristic-sensitive determinants is standard; no new
general Lefschetz mechanism is claimed.

The original primary question, pp.~1594--1596 of~\cite{romer}, and
the introduction, Theorem B, chain setup, and Theorem 9.2 of the
18 June 2025 arXiv version of~\cite{ss} were read for this comparison.
The latter paper's entire proof was not audited and none of its theorems
is needed in the proof here. Bounded literature and public-record searches
on 22 October 2026 authenticated no earlier publication of this exact
construction. This is not an exhaustive novelty or exclusive-priority
certificate, and no claim about the present open status of stronger
variants is made. The project's earlier proof is the unpublished
\href{https://github.com/AlecKriebel/Math/pull/1337}{PR1337},
head \texttt{aa737e6}; it is a source of this note, not an earlier
published solution. The internal catalog identifier is 30004534;
the original-question DOI is not a solution DOI.

The accompanying source and verification archive contains a new
standard-library exact checker. It independently constructs small
quadratic relation quotients by integral unit-pivot elimination and checks
multiplication ranks over finite fields, including exhaustive small
coefficient sets, injection compatibility, width-three generation, and
deliberately false controls. These finite checks supplement the universal
written proof; they do not certify priority or replace a proof for arbitrary
widths and fields. Source hashes, exact read scope, licenses, expected
outputs, and offline reproduction instructions are included.

\paragraph{Authorship and review disclosure.}
AI tools were used extensively in all phases: planning, mathematical
exploration, proof development and reconstruction, computation,
literature investigation, writing, editing, and adversarial verification.
Alec Kriebel is the named human author. This preprint is unrefereed,
with no conventional independent human peer review, journal acceptance,
or formal proof-assistant certification. Internal AI audits do not provide
those certifications. Authored prose is licensed CC BY 4.0; authored
verification and packaging code is MIT licensed. Third-party works
retain their own rights and are cited rather than redistributed.

\begin{thebibliography}{9}\small
\bibitem{romer} T.~R\"omer, \emph{Asymptotic properties of invariant chains},
in \emph{Lefschetz Properties in Algebra, Geometry and Combinatorics},
Oberwolfach Reports \textbf{17} (2020), report 31, 1594--1596.
\href{https://doi.org/10.4171/OWR/2020/31}{doi:10.4171/OWR/2020/31}.
\bibitem{ss} A.~Seceleanu and L.~\c Sega,
\emph{General symmetric ideals}, arXiv:2506.15459v1, 18 June 2025.
\href{https://arxiv.org/abs/2506.15459v1}{arxiv.org/abs/2506.15459v1}.
\bibitem{cook} D.~Cook II,
\emph{The Lefschetz properties of monomial complete intersections in positive characteristic},
J.~Algebra \textbf{369} (2012), 42--58.
\href{https://doi.org/10.1016/j.jalgebra.2012.07.015}{doi:10.1016/j.jalgebra.2012.07.015}.
\end{thebibliography}
\end{document}
